\section{The old whole affine blowup lift on its actual native fraction chart}
\begin{proposition}[Open pullback of the same original form lift]
\label{prop:whole-open-original-form-transfer}
\uses{thm:resume-original-lifted-form-open-naturality}
\lean{EnvelopeProperty.BlowupHomogeneousTautologicalLine.liftedFormSection_isIso_pullback_of_composite}
\leanok
Let $\mu:Z\to X$ be a scheme morphism, let $c:W\to Z$ be an open immersion,
and let $j:X\to\mathbf P_R^n$, $J$, and $F\in J_d$ be the original homogeneous
embedding, homogeneous ideal, and homogeneous form. Write $\beta_\mu(F)$ for
the original ideal-valued section, with target
$\mu^*j^*\mathcal O(d)\otimes(J\mathcal O_X\mathcal O_Z)$.
If $\beta_{c\mu}(F)$ is an isomorphism from $\mathcal O_W$ to its entire
ideal-tensor target, then $c^*\beta_\mu(F)$ is an isomorphism as well.
\end{proposition}
\begin{proof}
Put $M=j^*\mathcal O(d)$ and $D=J\mathcal O_X$. Since $c$ is open, pullback
of the whole ideal module is canonically the module of its whole comap ideal.
Together with the tensor pullback comparison and the canonical comparison
$c^*\mu^*M\simeq(c\mu)^*M$, this gives an isomorphism
$E:c^*(\mu^*M\otimes(D\mathcal O_Z))\simeq
 (c\mu)^*M\otimes(D\mathcal O_W)$.
The original-form naturality theorem, proved by comparing both included
sections and cancelling the monic ideal-tensor inclusion, gives
$(c^*\beta_\mu(F))E=\beta_{c\mu}(F)$.
The right side is an isomorphism by the hypothesis. Hence the composite on
the left is an isomorphism, and cancelling the isomorphism $E$ shows that
$c^*\beta_\mu(F)$ is an isomorphism. The target comparison and its section
identity have been derived; neither is a matching hypothesis.
\end{proof}
\begin{proposition}[The original affine blowup lift on its actual fraction chart]
\label{prop:whole-original-affine-blowup-fraction-basis}
\uses{prop:whole-open-original-form-transfer,thm:resume-original-form-actual-fraction-basis}
\lean{EnvelopeProperty.BlowupHomogeneousTautologicalLine.liftedFormSection_pullback_isIso_on_originalFractionChart}
\leanok
Let $j:X\to\mathbf P_R^n$, $J$, and $F\in J_d$ be the original projective
embedding, homogeneous ideal, and homogeneous form. Let $U$ be an affine open
of $X$ with integral coordinate ring, contained in the original $i$th projective
coordinate open. Put $D=J\mathcal O_X$, let $g\in D(U)$ be the original pulled
coefficient of $F$, and assume $g\ne0$. Put $A=\Gamma(U,\mathcal O_X)$,
$B=\operatorname{Bl}_{D(U)}(\operatorname{Spec}A)$,
$W=\operatorname{Spec}(A[D(U)/g])$, and let $c:W\to B$ be its actual native
fraction-chart inclusion. If $\mu:B\to X$ is the original affine blowup
structure followed by the actual affine-open map, then the pullback
$c^*\beta_\mu(F)$ of the original ideal-valued form lift is an isomorphism.
Its source is $\mathcal O_W$ and its target is the entire pullback of
$\mu^*j^*\mathcal O(d)\otimes(D\mathcal O_B)$.
\end{proposition}
\begin{proof}
The original ideal is generated by the pullbacks of all its original
homogeneous members. In particular the original coefficient $g$ lies in $D(U)$,
so the native fraction chart and its actual open inclusion $c$ are defined.
The proved original fraction-chart basis theorem says that the
original lift along $(cs)a$ is an isomorphism, where $s:B\to\operatorname{Spec}A$
is the actual affine blowup structure and $a:\operatorname{Spec}A\to X$ is the
actual affine-open map. Associativity identifies $(cs)a$ with $c\mu$, where
$\mu=sa$. Transporting the same ideal-valued section along that equality shows
that $\beta_{c\mu}(F)$ is an isomorphism. Apply
\ref{prop:whole-open-original-form-transfer}, using the actual native open
immersion $c$, to conclude that $c^*\beta_\mu(F)$ is an isomorphism.
All coefficient membership, canonical comparisons, and composite-basis facts
are derived from the original data. No whole-line, section-matching,
regularity, principal-ideal, or norm-matching certificate is an additional
hypothesis of this assertion.
\end{proof}
